\documentclass[11pt]{article}

\usepackage[T1]{fontenc}
\usepackage{lmodern}
\usepackage{microtype}
\usepackage[a4paper,margin=30mm]{geometry}
\usepackage{amsmath,amssymb,amsthm,mathtools}
\usepackage{booktabs}
\usepackage{enumitem}
\usepackage[hidelinks]{hyperref}
\hypersetup{
  pdftitle={The exponential rate for four-term-progression-free subset-sum sets},
  pdfauthor={Anonymous/deleted Reddit contributor; The Clankers},
  pdfsubject={Erdos Problem 817, k=4 exponential rate},
  pdfkeywords={subset sums, arithmetic progressions, Erdos Problem 817, Lean}
}

\newtheorem{theorem}{Theorem}[section]
\newtheorem{lemma}[theorem]{Lemma}
\newtheorem{proposition}[theorem]{Proposition}
\newtheorem{corollary}[theorem]{Corollary}
\theoremstyle{remark}
\newtheorem{remark}[theorem]{Remark}

\newcommand{\Z}{\mathbb Z}
\newcommand{\F}{\mathbb F}
\newcommand{\supp}{\operatorname{supp}}
\newcommand{\AP}{\mathrm{AP}}

\title{The exponential rate for four-term-progression-free subset-sum sets}
\author{Anonymous/deleted Reddit contributor\\
  \small proof construction\\[3pt]
  The Clankers\\
  \small independent reconstruction, verification, and formalization}
\date{13 September 2026}

\begin{document}
\maketitle

\begin{abstract}
For a finite set $A$ of positive integers, let $H(A)$ be its set of subset
sums, including zero, and let $g_4(n)$ be the least $N$ for which an
$n$-element set $A\subseteq\{1,\ldots,N\}$ has $H(A)$ free of nonconstant
four-term arithmetic progressions.  We give a complete reconstruction and
independent check of a proof posted by a now-deleted Reddit account.  The proof
determines the exponential rate:
\[
  \lim_{n\to\infty}g_4(n)^{1/n}=19^{1/3}.
\]
The lower bound classifies every relation with coefficients in
$\{-2,-1,0,1,2\}$ into disjoint minimal signed-relation blocks and counts the
resulting ternary collision classes.  The upper bound uses the block
$\{1,7,8\}$, whose seven subset-sum residues form a four-progression-free digit
set modulo $19$, and lifts it without carries in base $19$.  We state every
compressed step explicitly, record small typographical corrections to the
source, and supply Lean and executable certificates whose exact scopes are
listed in the final section.
\end{abstract}

\paragraph{Attribution and claim of credit.}
The construction and proof route in this note are attributed to the
anonymous contributor who posted the source on
r/LLMmathematics under an account that was subsequently deleted
\cite{AnonymousReddit2026}.  The linked ChatGPT conversation is also retained
as a public source locator \cite{ChatGPTShare2026}.  The Clankers do not
claim priority for that construction; ``The Clankers'' denotes the independent
reconstruction, checking, exposition, and formalization recorded here.

The moderators retain contemporaneous private mod-mail or notification
evidence of the original Reddit username.  The contributor is invited to
request any public attribution they prefer; the private record can be used to
validate a claimant without publishing the deleted username.  In short, the
proof is theirs, and this note is intended to make it checkable and return it
to them with its provenance intact.

\section{Problem and result}

For a finite set $A=\{a_1,\ldots,a_n\}$ of distinct positive integers, define
\[
  H(A)=\left\{\sum_{i=1}^n\varepsilon_i a_i:
       \varepsilon_i\in\{0,1\}\right\}.
\]
A four-term arithmetic progression is an ordered tuple
$(x,x+d,x+2d,x+3d)$; it is nonconstant when $d\ne0$.  Define $g_4(n)$ to be
the least positive integer $N$ for which there is an $n$-element set
$A\subseteq\{1,\ldots,N\}$ such that $H(A)$ contains no nonconstant four-term
arithmetic progression.

This is the $k=4$ exponential-rate subproblem of Erd\H{o}s Problem 817, which
originates in work of Erd\H{o}s and S{\'a}rk{\"o}zy
\cite{ErdosSarkozy1992,Bloom817}.  Korsky's recent general results give, for
fixed $k\ge4$, a lower exponential base $(k-1)/(k-2)$ and a separate
restricted-digit upper construction \cite[Theorems 1.2 and 1.3]{Korsky2026}.
The result below specializes to $k=4$ but obtains the exact exponential rate.
It does not address all of Problem 817.  In particular, it is logically
separate from the highlighted $k=3$ question, to which Costa has recently
given a negative answer by proving
$\liminf_{n\to\infty}g_3(n)/3^n=0$ \cite{Costa2026}.

\begin{theorem}[Main theorem]\label{thm:main}
For every positive integer $n$,
\[
  \frac{19^{n/3}-1}{2n}
  \le g_4(n)
  \le 8\,19^{\lceil n/3\rceil-1}.
\]
Consequently,
\[
  \lim_{n\to\infty}g_4(n)^{1/n}=19^{1/3}.
\]
\end{theorem}

\section{Relation splitting}

Assume throughout Sections~\ref{sec:split}--\ref{sec:lower} that
$A=\{a_1,\ldots,a_n\}\subseteq\{1,\ldots,N\}$ and that $H(A)$ is free of
nonconstant four-term arithmetic progressions.

\subsection{The two coefficient layers}\label{sec:split}

\begin{lemma}[Relation splitting]\label{lem:split}
Suppose
\[
  \sum_{i=1}^n c_i a_i=0,
  \qquad c_i\in\{-2,-1,0,1,2\}.
\]
Define
\[
  P_j=\sum_{c_i=j}a_i,
  \qquad N_j=\sum_{c_i=-j}a_i
  \qquad(j=1,2).
\]
Then $P_1=N_1$ and $P_2=N_2$.
\end{lemma}

\begin{proof}
The relation is
\[
  P_1+2P_2=N_1+2N_2.
\]
Set
\[
  X_0=P_1+N_2,\quad X_1=P_1+P_2,\quad
  X_2=N_1+N_2,\quad X_3=N_1+P_2.
\]
The four coefficient classes are disjoint, so every $X_j$ is a subset sum of
$A$.  Put $d=P_2-N_2$.  The displayed relation gives
$N_1-P_1=2d$, and direct subtraction, without changing coordinates, gives
\[
  X_1-X_0=d,
\]
\[
  X_2-X_1=(N_1-P_1)+(N_2-P_2)=2d-d=d,
\]
and
\[
  X_3-X_2=P_2-N_2=d.
\]
Thus $(X_0,X_1,X_2,X_3)$ is a four-term arithmetic progression in $H(A)$.
It must be constant, so $d=0$.  Hence $P_2=N_2$, and the original relation
then gives $P_1=N_1$.
\end{proof}

Equivalently, if a coefficient vector in $[-2,2]^n\cap\Z^n$ is a relation,
then its magnitude-one layer and its magnitude-two layer, after division by
two, are separately signed relations.

\subsection{Minimal signed-relation blocks}

A \emph{signed relation} is a nonzero vector
$r\in\{-1,0,1\}^n$ satisfying $\sum_i r_i a_i=0$.  It is \emph{minimal} if
no nonzero signed relation has support properly contained in $\supp r$.

\begin{lemma}[Minimal blocks are disjoint]\label{lem:disjoint}
If $r$ and $s$ are minimal signed relations, then either
$\supp r\cap\supp s=\varnothing$ or $s=\pm r$.
\end{lemma}

\begin{proof}
First let $r$ be minimal and let $s$ be any signed relation.  Define
\[
  E_+=\{i:r_i=s_i\ne0\},
  \qquad E_-=\{i:r_i=-s_i\ne0\}.
\]
Apply Lemma~\ref{lem:split} to $r+s$.  Its magnitude-two layer, divided by
two, is the restriction $r\mathbf 1_{E_+}$, so
\[
  \sum_{i\in E_+}r_i a_i=0.
\]
Apply the same lemma to $r-s$ to obtain
\[
  \sum_{i\in E_-}r_i a_i=0.
\]
If the supports of $r$ and $s$ meet, one of $E_+$ and $E_-$ is nonempty.
Minimality of $r$ forces that nonempty set to be all of $\supp r$.  Since
$E_+$ and $E_-$ are disjoint, exactly one is all of $\supp r$.  Hence $s=r$
on $\supp r$ or $s=-r$ there.

If $s$ is also minimal, this gives $\supp r\subseteq\supp s$.  A strict
inclusion would make $r$ a signed relation properly supported inside $s$,
contradicting minimality of $s$.  The supports are equal and $s=\pm r$.
\end{proof}

Choose one representative
\[
  r^{(1)},\ldots,r^{(h)}
\]
from every opposite pair of minimal signed relations, and write
$B_j=\supp r^{(j)}$.  The $B_j$ are pairwise disjoint.

\begin{lemma}[Every signed relation is a block sum]\label{lem:signed-decomp}
Every signed relation $q$ has a unique expression
\[
  q=\sum_{j=1}^h\alpha_j r^{(j)},
  \qquad \alpha_j\in\{-1,0,1\}.
\]
\end{lemma}

\begin{proof}
For existence, if $q\ne0$, choose a signed relation $u$ of inclusion-minimal
support among those supported in $\supp q$.  It is globally minimal.  The
first paragraph of the proof of Lemma~\ref{lem:disjoint} says that $q$ agrees
with either $u$ or $-u$ on all of $\supp u$.  Subtract that agreeing copy.
The remaining vector is zero or a signed relation, is zero on $\supp u$, and,
when nonzero, has strictly smaller support.  Induction on support size proves
existence.
Uniqueness follows coordinatewise because the block supports are disjoint and
every coordinate on a block is nonzero in its block vector.
\end{proof}

\begin{lemma}[Complete short-kernel classification]\label{lem:short-kernel}
If $c\in\{-2,-1,0,1,2\}^n$ and $\sum_i c_i a_i=0$, then uniquely
\[
  c=\sum_{j=1}^h t_j r^{(j)},
  \qquad t_j\in\{-2,-1,0,1,2\}.
\]
Conversely, every such block sum is a relation in
$\{-2,-1,0,1,2\}^n$.
\end{lemma}

\begin{proof}
Retain the two layers rather than merging them.  Define
\[
  u_i=\begin{cases}c_i,&|c_i|=1,\\0,&|c_i|\ne1,\end{cases}
  \qquad
  v_i=\begin{cases}c_i/2,&|c_i|=2,\\0,&|c_i|\ne2.\end{cases}
\]
Then $c=u+2v$ and $\supp u\cap\supp v=\varnothing$.  Lemma~\ref{lem:split}
says that each nonzero layer among $u$ and $v$ is a signed relation.  A zero
layer has the unique empty (all-zero) block decomposition.  Applying
Lemma~\ref{lem:signed-decomp} to each nonzero layer, and using that convention
for a zero layer, gives
\[
  u=\sum_j\alpha_jr^{(j)},
  \qquad v=\sum_j\beta_jr^{(j)},
  \qquad \alpha_j,\beta_j\in\{-1,0,1\}.
\]
If $\alpha_j\ne0$, the entire block $B_j$ lies in $\supp u$; if
$\beta_j\ne0$, the entire block lies in $\supp v$.  Their supports are
disjoint, so $\alpha_j\beta_j=0$ for every $j$.  Hence
\[
  t_j=\alpha_j+2\beta_j\in\{-2,-1,0,1,2\},
\]
which proves existence.  Disjoint block supports give uniqueness.  Conversely,
each block vector is a relation, and disjointness ensures that every coordinate
of $\sum_jt_jr^{(j)}$ remains in the stated coefficient set.
\end{proof}

Every block has size at least three.  A one-coordinate signed relation would
force some positive $a_i$ to be zero.  A two-coordinate signed relation would
either add two positive integers to zero or force two distinct elements of
$A$ to be equal.  Both are impossible.

\section{Counting ternary collision classes}\label{sec:lower}

Define
\[
  \Phi:\{0,1,2\}^n\longrightarrow\Z,
  \qquad \Phi(x)=\sum_{i=1}^n x_i a_i,
\]
and let $T(A)$ be its image.  By Lemma~\ref{lem:short-kernel},
\[
  \Phi(x)=\Phi(y)
  \quad\Longleftrightarrow\quad
  x-y=\sum_j t_jr^{(j)}
\]
for unique scalars $t_j\in\{-2,-1,0,1,2\}$.

Fix a block $B_j$ of size $m$.  On the block cube
$\{0,1,2\}^{B_j}$, reverse the coordinate on every position at which
$r^{(j)}_i=-1$ by the exact cube involution
\[
  (F_jx)_i=
  \begin{cases}
    x_i,&r^{(j)}_i=1,\\
    2-x_i,&r^{(j)}_i=-1.
  \end{cases}
\]
Then
\[
  x-y=t\,r^{(j)}|_{B_j}
  \quad\Longleftrightarrow\quad
  F_jx-F_jy=t(1,\ldots,1).
\]
Each diagonal class in $\{0,1,2\}^m$ has exactly one representative with
minimum coordinate zero: subtract the minimum coordinate.  There are $3^m$
ternary vectors, and exactly $2^m$ have all coordinates in $\{1,2\}$.  Thus
the block contributes exactly
\[
  3^m-2^m
\]
collision classes.  If $t$ coordinates lie outside the union of all blocks
and $m_j=|B_j|$, disjointness makes the quotient a direct product and gives
the exact count
\begin{equation}\label{eq:ternary-count}
  |T(A)|=3^t\prod_{j=1}^h(3^{m_j}-2^{m_j}),
  \qquad t+\sum_{j=1}^h m_j=n.
\end{equation}

For $m\ge3$, put $f_m=3^m-2^m$.  We have $f_3=19$ and
\[
  f_{m+1}=3f_m+2^m>3f_m.
\]
Since $3>19^{1/3}$, induction from $m=3$ yields
\[
  3^m-2^m\ge19^{m/3}\qquad(m\ge3).
\]
Also $3^t\ge19^{t/3}$.  Equation~\eqref{eq:ternary-count} therefore gives
\[
  |T(A)|\ge19^{n/3}.
\]
Every ternary sum is an integer in the interval
\[
  0\le\Phi(x)\le2\sum_i a_i\le2nN.
\]
Consequently
\[
  19^{n/3}\le |T(A)|\le2nN+1,
\]
and hence
\begin{equation}\label{eq:lower}
  g_4(n)\ge\frac{19^{n/3}-1}{2n}.
\end{equation}

\begin{remark}[Finite strengthening]\label{rem:finite-lower}
Because the generators are distinct,
\[
  \sum_i a_i\le nN-\frac{n(n-1)}2.
\]
The same proof gives the independent finite refinement
\[
  g_4(n)\ge
  \left\lceil\frac{19^{n/3}+n(n-1)-1}{2n}\right\rceil.
\]
This refinement is not needed for the exponential rate.
\end{remark}

\section{The seven-digit construction modulo 19}

Let
\[
  L(u,v,w)=u+7v+8w
\]
and let
\[
  D=L(\{0,1\}^3)=\{0,1,7,8,9,15,16\}\subset\Z/19\Z.
\]
Whenever its elements are used as base-$19$ digits or in integer equalities,
$D$ denotes these displayed canonical representatives in
$\{0,1,\ldots,18\}\subset\Z$; reduction modulo $19$ is stated explicitly.
The residue $8$ has the two binary representatives $(0,0,1)$ and $(1,1,0)$;
this collision is retained throughout the argument.

\begin{lemma}[Bounded modular kernel]\label{lem:mod-kernel}
For $u,v,w\in\{-2,-1,0,1,2\}$,
\[
  u+7v+8w\equiv0\pmod{19}
  \quad\Longleftrightarrow\quad
  (u,v,w)=t(1,1,-1)
\]
for a unique $t\in\{-2,-1,0,1,2\}$.
\end{lemma}

\begin{proof}
The reverse implication is immediate from $1+7-8=0$.  For the forward
implication put
\[
  p=u+w,\qquad q=v+w.
\]
Then $|p|,|q|\le4$ and $u+7v+8w=p+7q$.  Its absolute value is at most $32$,
so divisibility by $19$ forces $p+7q\in\{-19,0,19\}$.

If $p+7q=19$, the bounds $|p|,|q|\le4$ force $(p,q)=(-2,3)$.  But
$v+w=3$ forces $w\ge1$, so $u+w\ge-1$, contradicting $p=-2$.  Negation
excludes $-19$.  Thus $p+7q=0$.  The bound $|p|\le4$ forces $q=0$, then
$p=0$, and therefore $u=v=-w$.
\end{proof}

\begin{proposition}[Digit rigidity]\label{prop:digit}
The set $D$ contains no nonconstant four-term arithmetic progression modulo
$19$.
\end{proposition}

\begin{proof}
Suppose $z_0,z_1,z_2,z_3\in D$ have vanishing consecutive second differences
modulo $19$.  Choose arbitrary binary representatives
$x_i\in\{0,1\}^3$ with $L(x_i)=z_i$.  Both vectors
\[
  x_0-2x_1+x_2,
  \qquad x_1-2x_2+x_3
\]
belong to $\{-2,-1,0,1,2\}^3$ and lie in the modular kernel of $L$.
Lemma~\ref{lem:mod-kernel} makes each a scalar multiple of $(1,1,-1)$.

Now retain the exact projection
\[
  \pi:\Z^3\longrightarrow\Z^2,
  \qquad \pi(a,b,c)=(a+c,b+c).
\]
Its kernel contains $(1,1,-1)$, and
\[
  L(a,b,c)=\pi_1(a,b,c)+7\pi_2(a,b,c).
\]
Applying $\pi$ shows that
$\pi(x_0),\pi(x_1),\pi(x_2),\pi(x_3)$ is an ordinary integer four-term
progression in $\{0,1,2\}^2$.  In either coordinate its first and fourth
values differ by three times the integer common difference, but both lie in
$[0,2]$.  The common difference is therefore zero in each coordinate.  All
four projections, and hence all four $L$-values, are equal.  Thus the original
modular progression is constant.  Either representative of the duplicated
digit $8$ gives the same projection $(1,1)$, so no representative choice has
been suppressed.
\end{proof}

\section{Carry-free base-19 lifting}

For $m\ge1$, define the set of $3m$ distinct positive generators
\[
  A_m=\bigcup_{j=0}^{m-1}19^j\{1,7,8\}.
\]
Each subset chooses a binary triple in each block.  Hence its contribution at
base-19 position $j$ is a digit in $D$.  Conversely, every digit of $D$ is a
subset sum of $\{1,7,8\}$.  Since the largest digit is $16<19$, no carry
occurs, and therefore
\begin{equation}\label{eq:language}
  H(A_m)=\left\{\sum_{j=0}^{m-1}d_j19^j:d_j\in D\right\}.
\end{equation}

\begin{proposition}[Least-differing-digit lift]\label{prop:lift}
For every $m\ge1$, the set $H(A_m)$ contains no nonconstant four-term
arithmetic progression over the integers.
\end{proposition}

\begin{proof}
Suppose
\[
  y_i=y_0+i\delta\in H(A_m)\qquad(0\le i\le3)
\]
for a nonzero integer $\delta$.  Let $r=v_{19}(|\delta|)$, so
$\delta=19^r e$ with $19\nmid e$.  The exponent satisfies $r<m$: indeed,
every element of $H(A_m)$ lies in
\[
  \left[0,16\sum_{j=0}^{m-1}19^j\right]
  =\left[0,\frac89(19^m-1)\right]\subset[0,19^m),
\]
so a nonzero difference of two such elements cannot be divisible by $19^m$.

The four numbers are congruent modulo $19^r$.  By uniqueness of the standard
base-19 expansion in \eqref{eq:language}, their lower $r$ digits are identical.
Let $c$ be the common integer represented by those digits and put
\[
  z_i=\frac{y_i-c}{19^r}.
\]
Each $z_i$ still has all base-19 digits in $D$, and the four $z_i$ form an
integer progression of common difference $e$.  Reducing modulo $19$ now reads
their least remaining digits and produces a four-term progression in $D$ with
nonzero common difference $e\bmod19$.  This contradicts
Proposition~\ref{prop:digit}.
\end{proof}

The generators in $A_m$ are distinct because their $19$-adic valuations give
the block index and the multipliers $1,7,8$ are distinct within a block.  Thus
\[
  |A_m|=3m,
  \qquad \max A_m=8\,19^{m-1},
\]
and Proposition~\ref{prop:lift} gives
\[
  g_4(3m)\le8\,19^{m-1}.
\]
For arbitrary $n\ge1$, take $m=\lceil n/3\rceil$ and delete any $3m-n$
generators.  If $B\subseteq A_m$, then $H(B)\subseteq H(A_m)$; deletion cannot
create a progression.  Therefore
\begin{equation}\label{eq:upper}
  g_4(n)\le8\,19^{\lceil n/3\rceil-1}.
\end{equation}

\section{Completion of the proof}

Equations~\eqref{eq:lower} and \eqref{eq:upper} prove the finite bounds in
Theorem~\ref{thm:main}.  For the asymptotic statement, the $n$th root of the
lower expression is
\[
  19^{1/3}\bigl(1-19^{-n/3}\bigr)^{1/n}(2n)^{-1/n},
\]
which tends to $19^{1/3}$, while the $n$th root of the upper expression is
\[
  19^{\lceil n/3\rceil/n}(8/19)^{1/n},
\]
which has the same limit.  The squeeze theorem proves
\[
  \lim_{n\to\infty}g_4(n)^{1/n}=19^{1/3}.
\]

\section{Corrections and exact verification status}

The source post loses superscripts in two plain-text phrases: ``$2m$ vectors''
must be $2^m$ vectors, and ``$19r$'' must be $19^r$.  It also compresses the
disjoint-layer sentence in Lemma~\ref{lem:short-kernel}, the common-low-digit
removal in Proposition~\ref{prop:lift}, and the fact that $r<m$.  Those steps
are supplied above.  Finally, the statement that a three-generator block has
nineteen ternary classes requires the ambient four-progression-free hypothesis
and the complete short-kernel classification.  It is not a statement about an
arbitrary triple with one visible signed relation.

The accompanying Lean 4 modules certify, without \texttt{sorry}, the following exact
claims:
\begin{enumerate}[label=(\roman*)]
  \item the bounded kernel in Lemma~\ref{lem:mod-kernel};
  \item the two-second-difference rigidity of the seven-digit set;
  \item the recursive base-19 lift for every word length;
  \item the identification of each binary block choice with one of the seven
        digits;
  \item the relation-splitting algebra in Lemma~\ref{lem:split};
  \item the distinctness and positivity of the $3m$ generators, their exact
        subset-sum representation by binary block choices, and the bound
        $\max A_m\le8\,19^{m-1}$ for $m\ge1$;
  \item the passage to an $n$-element subcollection, proving the finite upper
        bound \eqref{eq:upper} for every positive integer $n$.
\end{enumerate}
The extension from the digit language to the finite-generator upper bound was
formalized by \href{https://github.com/sneed-and-feed}{\texttt{sneed-and-feed}}
in \href{https://github.com/KokunoYumeto/erdos-problem-817-workbench/pull/1}{pull request~1}.
The construction and mathematical proof remain those of the original
anonymous contributor.  Both modules were independently elaborated with
Lean~4.32.0 and the pinned Mathlib dependency; the axiom audit of all $44$
theorem declarations in the extension reports only \texttt{propext},
\texttt{Classical.choice}, and \texttt{Quot.sound}, or subsets thereof.

The Lean modules do not yet certify the full minimal-block decomposition, the
global product cardinality in \eqref{eq:ternary-count}, or the real-asymptotic
squeeze.  Those parts are proved in ordinary mathematics in this note and are
not mislabeled as Lean-complete.

The executable certificate independently enumerates all $125$ bounded
coefficient triples, all $342$ start/nonzero-step modular progressions, all
seven-digit second-difference quadruples, the lifted languages through four
blocks, and the exact deletion choices through nine generators.  A separate
stress test checks the minimal-block and ternary-count statements over a stated
finite search domain.  Finite enumeration is corroborating evidence, not a
replacement for the general proofs.

\paragraph{Novelty status.}
The checked sources below do not state the exact $19^{1/3}$ rate, but this is
not a priority claim.  A bounded literature search cannot establish novelty.
The broad catalogue problem asks for estimates across $k\ge3$
\cite{Bloom817}.  Costa's negative answer for the historical $k=3$ lower-bound
question and the present $k=4$ exponential-rate theorem concern different
subproblems and do not exhaust that broader request.  The present theorem may
be new, overlooked, independently rediscovered, or in need of further
community review.

\begingroup
\small
\bibliographystyle{plain}
\bibliography{references}
\endgroup

\end{document}
